% Part: first-order-logic
% Chapter: natural-deduction
% Section: propositional-rules

\documentclass[../../../include/open-logic-section]{subfiles}

\begin{document}

\iftag{FOL}
      {\olfileid[fr]{fol}{ntd}{prl}}
      {\olfileid[fr]{pl}{ntd}{prl}}

\olsection{Règles propositionnelles}

\Needspace{11\baselineskip}
\subsection{Règles pour $\land$}

\begin{defish}
\AxiomC{$!A$}
\AxiomC{$!B$}
\RightLabel{\Intro{\land}}
\BinaryInfC{$!A \land !B$}
\DisplayProof
\hfill
\begin{tabular}{r}
\AxiomC{$!A \land !B$}
\RightLabel{\Elim{\land}}
\UnaryInfC{$!A$}
\DisplayProof
\\[3ex]
\AxiomC{$!A \land !B$}
\RightLabel{\Elim{\land}}
\UnaryInfC{$!B$}
\DisplayProof
\end{tabular}
\end{defish}

\Needspace{11\baselineskip}
\subsection{Règles pour $\lor$}

\begin{defish}
\begin{tabular}{r}
\AxiomC{$!A$}
\RightLabel{\Intro{\lor}}
\UnaryInfC{$!A \lor !B$}
\DisplayProof
\\[3ex]
\AxiomC{$!B$}
\RightLabel{\Intro{\lor}}
\UnaryInfC{$!A \lor !B$}
\DisplayProof
\end{tabular}
\hfill
\AxiomC{$!A \lor !B$}
\AxiomC{$\Discharge{!A}{n}$}
\DeduceC{$!C$}
\AxiomC{$\Discharge{!B}{n}$}
\DeduceC{$!C$}
\DischargeRule{\Elim{\lor}}{n}
\TrinaryInfC{$!C$}
\DisplayProof
\end{defish}

\Needspace{11\baselineskip}
\subsection{Règles pour $\lif$}

\begin{defish}
\AxiomC{$\Discharge{!A}{n}$}
\DeduceC{$!B$}
\DischargeRule{\Intro{\lif}}{n}
\UnaryInfC{$!A \lif !B$}
\DisplayProof
\hfill
\AxiomC{$!A \lif !B$}
\AxiomC{$!A$}
\RightLabel{\Elim{\lif}}
\BinaryInfC{$!B$}
\DisplayProof
\end{defish}

\Needspace{11\baselineskip}
\subsection{Règles pour $\lnot$}

\begin{defish}
\AxiomC{$\Discharge{!A}{n}$}
\noLine
\DeduceC{$\lfalse$}
\DischargeRule{\Intro{\lnot}}{n}
\UnaryInfC{$\lnot !A$}
\DisplayProof
\hfill
\AxiomC{$\lnot !A$}
\AxiomC{$!A$}
\RightLabel{\Elim{\lnot}}
\BinaryInfC{$\lfalse$}
\DisplayProof
\end{defish}

\Needspace{11\baselineskip}
\subsection{Règles pour $\lfalse$}

\begin{defish}
\AxiomC{$\lfalse$}
\RightLabel{\FalseInt}
\UnaryInfC{$!A$}
\DisplayProof
\hfill
\AxiomC{$\Discharge{\lnot !A}{n}$}
\DeduceC{$\lfalse$}
\DischargeRule{\FalseCl}{n}
\UnaryInfC{$!A$}
\DisplayProof
\end{defish}

Les règles $\Intro{\lnot}$ et $\FalseCl$ sont très proches : la
première permet de dériver l'!!{sentence} nié~$\lnot !A$, tandis que
la seconde permet de dériver~$!A$, sans lui ajouter de négation.

Lorsqu'une règle indique que l'on peut décharger une hypothèse,
il s'agit d'une possibilité, et non d'une obligation. Par exemple,
la règle $\Intro{\lif}$ permet de décharger un nombre quelconque
d'hypothèses de la forme~$!A$ dans la !!{derivation} de la
prémisse~$!B$, y compris zéro.

\begin{editorial}
La source qualifie ici $!A$ de « positif ». Ce mot signifie seulement
que la règle n'ajoute pas de négation devant $!A$ : le schéma n'exige
pas que $!A$ soit lui-même dépourvu de négation.
\end{editorial}

\end{document}
